\section{Conclusions and Future Work}

In this study, we conduct systematic investigation into why autonomous
research agents fail. This investigation results in \textbf{\bench}: a
collection of 800 research trajectories with artifact-aware, process-level
failure annotations, elicited from eight harness--model combinations on a
100-task, seven-domain suite spanning the full scientific lifecycle. To enable
\bench's systematic annotation and analysis, we develop \textbf{ARFT},
the first \add{systematic }%
failure taxonomy for autonomous research agents, organizing 45
failure patterns on two cross-cutting axes: the lifecycle stage at which a
failure manifests and its underlying root cause. For scalable annotation we
develop an agent-as-a-judge annotator, validated against human expert annotation on a 50-trajectory sample (\S\ref{sec:judge:validation}). Our analysis reveals that all three cognitive root causes converge on a single
overarching limitation: current agents lack a \textbf{metacognitive
loop}---the ability to check what they produced against what they found, act
when it does not hold up, and question whether the path they took was sound.

We are excited about the potential of autonomous research agents, but their
adoption in science hinges on reliability that outcome-only benchmarks cannot
certify. Our work contributes toward this goal through the public release of
\bench, ARFT, the task suite, and our annotator. \bench offers a
rich empirical basis for understanding current failure dynamics, while ARFT
provides a standardized language to diagnose, attribute, and mitigate these
failures. Several directions follow naturally from our findings. First, our
three insights suggest that artifact-aware evaluation---scoring the run
directory alongside the final report---could catch failures that endpoint
scoring systematically misses, and developing such evaluation protocols is a
natural next step. Second, the concentration of failures in the
self-verification stage (F.1--F.4 together account for 14.1\% of all hits, and F.4 alone appears in 82.5\% of analyses) suggests that the review stage is where interventions may yield the
highest return, and studying how different review architectures affect failure
rates is a promising direction. Third, while the core failure patterns are
shared across all eight systems, fabrication rates vary substantially,
indicating that system-specific mitigation guided by the per-model breakdowns
in our appendix is underexplored. More broadly, by tracing
all three cognitive root causes to a shared metacognitive deficit, our
analysis identifies the metacognitive loop as a central target for next-generation language models and agentic systems working toward genuine scientific discovery.
